
% COLONNA X COMPATIBILE CON LONGTABLE
\newcolumntype{Y}{>{\raggedright\arraybackslash}p{8cm}}

\section{Requisiti funzionali e non funzionali}

\subsection{Requisiti funzionali}

\begin{longtable}{@{}p{1.2cm}p{4.5cm}Y@{}}

\caption{Requisiti funzionali.}\\

\toprule
\textbf{ID} & \textbf{Nome} & \textbf{Descrizione} \\
\midrule
\endfirsthead

\toprule
\textbf{ID} & \textbf{Nome} & \textbf{Descrizione} \\
\midrule
\endhead

\FRid{1} & Ricerca per signature &
Il sistema deve consentire di cercare un testo per
\code{track\_name} + \code{artist\_name}, con eventuale filtro su
\code{album\_name} e \code{duration}. Realizzato da
\endpoint{GET}{/api/get} e \endpoint{GET}{/api/get-cached}. \\

\FRid{2} & Recupero per ID &
Deve essere possibile recuperare un testo conoscendo il suo identificativo
univoco. Realizzato da \endpoint{GET}{/api/get/<id>}. \\

\FRid{3} & Ricerca generica &
L'utente deve poter cercare testi per parola chiave (\code{q}) o per
combinazione di titolo, artista, album. Realizzato da
\endpoint{GET}{/api/search}. \\

\FRid{4} & Pubblicazione anonima &
Un utente non registrato deve poter pubblicare un nuovo testo dopo aver
superato una sfida \emph{Proof of Work}. Realizzato dalla coppia
\endpoint{POST}{/api/request-challenge} +
\endpoint{POST}{/api/publish} con header \code{X-Publish-Token}. \\

\FRid{5} & Pubblicazione autenticata &
Un utente autenticato (Bearer token) deve poter pubblicare un nuovo testo
senza risolvere la PoW. Il brano viene attribuito al suo username nel
campo \code{submittedBy}. \\

\FRid{6} & Registrazione utente &
Deve essere possibile creare un account fornendo username, email e
password. Realizzato da \endpoint{POST}{/auth/register}. \\

\FRid{7} & Login utente &
L'utente deve potersi autenticare con username o email e password,
ricevendo un access token JWT. Realizzato da
\endpoint{POST}{/auth/login}. \\

\FRid{8} & Dati dell'utente corrente &
L'utente autenticato deve poter recuperare i propri dati di profilo.
Realizzato da \endpoint{GET}{/auth/me}. \\

\FRid{9} & Elenco dei propri brani &
L'utente autenticato deve poter ottenere la lista dei testi da lui
pubblicati. Realizzato da \endpoint{GET}{/api/me/songs}. \\

\FRid{10} & Modifica di un proprio brano &
L'utente autenticato deve poter modificare un brano di cui è autore.
Realizzato da \endpoint{PUT}{/api/songs/<id>}. \\

\FRid{11} & Cancellazione di un proprio brano &
L'utente autenticato deve poter cancellare un brano di cui è autore.
Realizzato da \endpoint{DELETE}{/api/songs/<id>}. \\

\FRid{12} & Pannello amministratore &
Un utente con flag \code{is\_admin} deve poter consultare l'intero
catalogo (inclusi i brani anonimi), con filtri per query e per
autore. Realizzato da \endpoint{GET}{/admin/songs}
(richiede JWT + \code{is\_admin}). \\

\FRid{13} & Esplorazione del catalogo &
Il sistema deve fornire un elenco paginato dei brani, ordinabile per
data, titolo o artista. Realizzato da \endpoint{GET}{/api/explore}. \\

\FRid{14} & Lettura del testo sincronizzato &
L'applicazione mobile deve poter mostrare il testo sincronizzato a tempo
con la riproduzione, evidenziando la riga corrente. \\

\FRid{15} & Health check &
Deve essere disponibile un endpoint per verificare lo stato del backend
e visualizzarlo in home (utile per lo sviluppo).
Realizzato da \endpoint{GET}{/api/health}. \\

\bottomrule
\end{longtable}

\subsection{Requisiti non funzionali}

I requisiti non funzionali descrivono le proprietà qualitative del sistema,
i vincoli architetturali e i requisiti di sicurezza, affidabilità,
performance e manutenibilità che il sistema Stopify deve rispettare.

\begin{longtable}{@{}p{1.2cm}p{3cm}p{8.5cm}@{}}

\caption{Requisiti non funzionali.}\\

\toprule
\textbf{ID} & \textbf{Categoria} & \textbf{Descrizione} \\
\midrule
\endfirsthead

\toprule
\textbf{ID} & \textbf{Categoria} & \textbf{Descrizione} \\
\midrule
\endhead

\NFRid{1} & Sicurezza &
Le password degli utenti devono essere memorizzate esclusivamente tramite
hashing PBKDF2-SHA256 utilizzando le funzioni di
\texttt{werkzeug.security}. \\

\NFRid{2} & Sicurezza &
Gli endpoint autenticati devono richiedere un JSON Web Token (JWT)
valido e non scaduto. \\

\NFRid{3} & Sicurezza &
Le pubblicazioni anonime devono essere consentite solo previa
risoluzione corretta di una sfida crittografica Proof of Work (PoW). \\

\NFRid{4} & Sicurezza &
Il backend deve limitare la dimensione massima dei payload HTTP per
prevenire attacchi DoS tramite richieste eccessivamente grandi. \\

\NFRid{5} & Privacy &
Le password non devono mai essere memorizzate o trasmesse in chiaro. \\

\NFRid{6} & Privacy &
Il sistema non deve associare dati personali identificativi alle
pubblicazioni anonime effettuate tramite PoW. \\

\NFRid{7} & Performance &
Le operazioni di ricerca e recupero testi devono garantire tempi di
risposta compatibili con l’utilizzo su dispositivi mobili. \\

\NFRid{8} & Affidabilità &
Il backend deve restituire codici di stato HTTP coerenti con l’esito
delle operazioni richieste. \\

\NFRid{9} & Compatibilità &
Lo scambio dati tra client mobile e backend deve avvenire esclusivamente
tramite JSON codificato UTF-8 su protocollo HTTP/REST. \\

\NFRid{10} & Portabilità &
L’applicazione client deve poter essere eseguita su piattaforme Android,
iOS e Web tramite React Native ed Expo. \\

\NFRid{11} & Disponibilità &
Il backend deve esporre un endpoint di health check per verificare lo
stato del servizio. \\

\NFRid{12} & Manutenibilità &
Il backend deve essere organizzato secondo un’architettura modulare
suddivisa in routes, services e models. \\

\NFRid{13} & Scalabilità &
L’architettura del sistema deve consentire la sostituzione futura del
database SQLite con DBMS più scalabili senza modifiche sostanziali alle
API REST. \\

\NFRid{14} & Sicurezza &
In ambiente di produzione le comunicazioni devono avvenire tramite HTTPS
con TLS terminato da un reverse proxy. \\

\bottomrule
\end{longtable}